\section{Anneaux japonais}
\label{section:0.23-fr}

Les résultats de ce paragraphe seront complétés dans IV, (7.6) et (7.7).

\subsection{Anneaux japonais}
\label{subsection:0.23.1-fr}

\begin{definition}[23.1.1]
\label{0.23.1.1-fr}
On dit qu'un anneau intègre $A$ est un anneau japonais si, pour toute
extension finie $K'$ de son corps des fractions $K$, la fermeture intégrale
$A'$ de $A$ dans $K'$ est

\oldpage[0\textsubscript{IV-1}]{214}

un $A$-module de type fini (autrement dit, une $A$-algèbre finie). On dit
qu'un anneau $A$ est universellement japonais si toute $A$-algèbre intègre de
type fini est un anneau japonais.
\end{definition}

Il est clair que si $A$ est un anneau japonais, tout anneau de fractions
$S^{-1}A$ est alors un anneau japonais, car (avec les notations précédentes)
$S^{-1}A'$ est la fermeture intégrale de $S^{-1}A$ dans $K'$, et est
évidemment un $S^{-1}A$-module de type fini.

Un anneau intègre noethérien, même un \emph{anneau de valuation discrète},
n'est pas nécessairement un anneau japonais [30].

\begin{proposition}[23.1.2]
\label{0.23.1.2-fr}
Soient $A$ un anneau noethérien intègre, $K$ son corps des fractions. Si, pour
toute extension radicielle finie $K'$ de $K$, la fermeture intégrale de $A$
dans $K'$ est un $A$-module de type fini, alors $A$ est un anneau japonais.
\end{proposition}

\begin{proof}
Comme $A$ est noethérien, il suffit, pour vérifier que $A$ est un anneau
japonais, de voir que pour toute extension quasi-galoisienne finie $L$ de
$K$, la fermeture intégrale $B$ de $A$ dans $L$ est un $A$-module de type
fini. Or $L$ est une extension galoisienne de la plus grande extension
radicielle $K'$ de $K$ contenue dans $L$~; et si $A'$ est la fermeture
intégrale de $A$ dans $K'$, $B$ est la fermeture intégrale de $A'$ dans $L$.
Mais on sait que dans une extension séparable de $K'$, la fermeture intégrale
de $A'$ est un $A'$-module de type fini (Bourbaki, \emph{Alg. comm.},
chap.~V, §~1, n\textsuperscript{o}~6, cor.~1 de la prop.~20), d'où la
proposition.
\end{proof}

Il résulte de (23.1.2) que lorsque $K$ est de caractéristique $0$, dire que
$A$ est un anneau japonais signifie que sa \emph{clôture intégrale} $A'$ est
un $A$-module de type fini.

\begin{theorem}[23.1.3]
\label{0.23.1.3-fr}
\textnormal{(Tate).} Soient $A$ un anneau intègre noethérien, $x\neq0$ un
élément de $A$. Supposons vérifiées les conditions suivantes~:
\begin{enumerate}
  \item[(i)] $A$ est intégralement clos.
  \item[(ii)] $\mathfrak p=xA$ est premier et $A$ est séparé et complet pour
  la topologie $\mathfrak p$-préadique.
  \item[(iii)] $A/xA$ est un anneau japonais.
\end{enumerate}
Alors $A$ est un anneau japonais.
\end{theorem}

Nous utiliserons le lemme suivant~:

\begin{lemma}[23.1.3.1]
\label{0.23.1.3.1-fr}
Soient $A$ un anneau, $x$ un élément de $A$ non diviseur de $0$ dans $A$ et
tel que $xA=\mathfrak p$ soit premier~; alors, pour tout entier $n>0$,
l'image réciproque de $x^nA_{\mathfrak p}$ dans $A$ est $x^nA$.
\end{lemma}

\begin{proof}
En effet, supposons que $b\in A$ soit un élément tel que
$b/1=x^na/s$ dans $A_{\mathfrak p}$, où $a\in A$ et
$s\notin\mathfrak p$~; il existe donc $s'\notin\mathfrak p$ tel que
$s'sb=x^nas'$, d'où $s'sb\in\mathfrak p$, et comme
$s's\notin\mathfrak p$, cela entraîne $b\in\mathfrak p$, autrement dit
$b=xb'$ avec $b'\in A$~; puisque $x$ est non diviseur de $0$, on en conclut
$s'sb'=x^{n-1}as$ et il suffit de raisonner par récurrence sur $n$.
\end{proof}

\begin{proof}
Pour démontrer le théorème, on peut se borner au cas où le corps des fractions
$K$ de $A$ est de caractéristique $p>0$, puisque $A$ est intégralement clos.
Soit $K'$ une extension radicielle finie de $K$, de sorte qu'il existe une
puissance $q=p^e$ telle que $K'^q\subset K$~; en remplaçant $K'$ par une
extension radicielle plus grande, on peut même supposer qu'il existe $y\in K'$
tel que $y^q=x$. Soit $A'$ la fermeture intégrale de $A$ dans $K'$~; comme
$A$ est intégralement clos, $A'$ est l'ensemble des $x'\in K'$ tels que
$x'^q\in A'\cap K=A$. Posons $V=A_{\mathfrak p}$~; l'idéal maximal
$\mathfrak m=xV$ de $V$ étant principal, on sait que l'anneau local
noethérien $V$ est un \emph{anneau de valuation discrète} (17.1.4)~; comme
$V$ est intégralement clos, le même

\oldpage[0\textsubscript{IV-1}]{215}

raisonnement que ci-dessus montre que la fermeture intégrale $V'$ de $V$ dans
$K'$ est l'ensemble des $x'\in K'$ tels que $x'^q\in V$~; on sait
(Bourbaki, \emph{Alg. comm.}, chap.~VI, §~8, n\textsuperscript{o}~6,
prop.~6, et chap.~V, §~2, n\textsuperscript{o}~3, lemme~4) que $V'$ est un
\emph{anneau de valuation discrète}, dont l'idéal maximal $\mathfrak m'$ est
l'ensemble des $x'\in K'$ tels que $x'^q\in\mathfrak m$ et en outre le corps
résiduel $V'/\mathfrak m'$ est une extension \emph{finie} de
$V/\mathfrak m$ (\emph{loc. cit.}, chap.~VI, §~8,
n\textsuperscript{o}~1, lemme~2). Montrons que, pour tout entier $n>0$, on a
\begin{equation}
\label{0.23.1.3.2-fr}
  \mathfrak m'^n\cap A'=y^nA'.
  \tag{23.1.3.2}
\end{equation}

En effet, comme $y^q=x$, on a $y\in\mathfrak m'$, donc
$y^nA'\subset\mathfrak m'^n\cap A'$. Inversement, soit
$x'\in\mathfrak m'^n\cap A'$, et posons $x'=z'y^n$ avec $z'\in K'$. On a
$x'^q\in A'$~; d'autre part, $x'$ est somme de produits
$t'_1\cdots t'_n$ avec $t'_i\in\mathfrak m'$, donc
$t_i'^q\in\mathfrak m$, et l'on en conclut que
$x'^q\in\mathfrak m^n\cap A$. Or le lemme (23.1.3.1) prouve que
$\mathfrak m^n\cap A=x^nA$, donc on a $z'^qx^n\in x^nA$, ou encore
$z'^q\in A$ puisque $x\neq0$, ce qui établit (23.1.3.2).

Montrons en second lieu que sur $A'$ la topologie $xA'$-préadique est
\emph{séparée}~; cette topologie est aussi la topologie $yA'$-préadique de
$A'$, et il résulte de (23.1.3.2) que cette topologie est induite par la
topologie $\mathfrak m'$-préadique de $V'$ qui est séparée puisque $V'$ est
un anneau de valuation discrète.

Prouvons ensuite que $A'/xA'$ est un $A$-module de type fini. Comme $x=y^q$
et que $y^kA'/y^{k+1}A'$ est isomorphe à $A'/yA'$, on peut se borner à
montrer que $A'/yA'$ est un $A$-module de type fini~; mais la formule
(23.1.3.2) appliquée pour $n=1$ montre que $A'/yA'$ est un sous-anneau de
$V'/\mathfrak m'$, c'est-à-dire du corps résiduel de $V'$, qui est une
extension finie du corps résiduel $V/\mathfrak m$ de $V$. Or
$V/\mathfrak m$ est le corps des fractions de $A/\mathfrak p$, et comme
$A'$ est entier sur $A$, $A'/yA'$ est entier sur $A/\mathfrak p$, donc
contenu dans la fermeture intégrale de $A/\mathfrak p$ dans
$V'/\mathfrak m'$~; comme par hypothèse $A/\mathfrak p$ est un anneau
japonais noethérien, $A'/yA'$ est un $(A/\mathfrak p)$-module de type fini,
et \emph{a fortiori} un $A$-module de type fini.

Cela étant, la topologie $xA'$-préadique de $A'$ étant séparée, $A'$ est un
sous-anneau de son complété $\widehat{A'}$ pour cette topologie~; mais puisque
$A'/xA'$ est un $A$-module de type fini et que $A$ est séparé et complet pour
la topologie $xA$-préadique, $\widehat{A'}$ est un $A$-module de type fini en
vertu de (0\textsubscript{I}, 7.2.9), donc il en est de même de $A'$, C.Q.F.D.
\end{proof}

\begin{corollary}[23.1.4]
\label{0.23.1.4-fr}
Soit $A$ un anneau japonais noethérien et intégralement clos. Alors tout
anneau de séries formelles $A[[T_1,\ldots,T_r]]$ est un anneau japonais.
\end{corollary}

\begin{proof}
On sait que pour tout anneau noethérien intégralement clos $B$, l'anneau de
séries formelles $B[[T]]$ est noethérien et intégralement clos (Bourbaki,
\emph{Alg. comm.}, chap.~V, §~1, n\textsuperscript{o}~4, prop.~14)~; on peut
donc, par récurrence sur $n$, se borner à prouver que $A[[T]]$ est un anneau
japonais. Or l'élément $x=T$ vérifie toutes les conditions de (23.1.3), d'où
la conclusion.
\end{proof}

\begin{theorem}[23.1.5]
\label{0.23.1.5-fr}
\textnormal{(Nagata).} Tout anneau local noethérien intègre complet $A$ est
un anneau japonais.
\end{theorem}

\begin{proof}
On sait (19.8.8, (ii)) qu'il existe un sous-anneau $B$ de $A$ qui est un
anneau local complet et régulier, tel que $A$ soit une $B$-algèbre finie~; il
suffit évidemment de prouver

\oldpage[0\textsubscript{IV-1}]{216}

que $B$ est un anneau japonais, autrement dit, on peut se borner au cas où
$A$ est en outre \emph{régulier}. Raisonnons par récurrence sur
$n=\dim(A)$, le théorème étant évident pour $n=0$. Supposons donc $n>0$, et,
en désignant par $\mathfrak m$ l'idéal maximal de $A$, soit $x$ un élément de
$\mathfrak m-\mathfrak m^2$~; on sait (17.1.8) que $A/xA$ est un anneau
régulier, et en outre ((17.1.7) et (16.3.4)) que
$\dim(A/xA)=n-1$~; de plus, l'idéal $xA$ est fermé dans $A$
(0\textsubscript{I}, 7.3.5), donc $A/xA$ est complet. En vertu de l'hypothèse
de récurrence, $A/xA$ est donc un anneau japonais, et il résulte alors de
(23.1.3) qu'il en est de même de $A$.
\end{proof}

\begin{corollary}[23.1.6]
\label{0.23.1.6-fr}
Soient $A$ un anneau local noethérien intègre complet, $K$ son corps des
fractions, $K'$ une extension finie de $K$~; alors la fermeture intégrale
$A'$ de $A$ dans $K'$ est un anneau local complet.
\end{corollary}

\begin{proof}
On sait déjà par (23.1.5) que $A'$ est une $A$-algèbre finie, donc un anneau
semi-local noethérien complet, et par suite un \emph{composé direct}
d'anneaux locaux~; mais comme $A'$ est intègre, c'est nécessairement un
anneau local.
\end{proof}

\begin{proposition}[23.1.7]
\label{0.23.1.7-fr}
Soient $A$ un anneau local noethérien intègre, $K$ son corps des fractions,
$K'$ une extension finie de $K$, $A'$ la fermeture intégrale de $A$ dans
$K'$. On suppose que le complété $\widehat A$ est réduit, et l'on désigne par
$R$ son anneau total des fractions.
\begin{enumerate}
  \item[(i)] Si l'anneau $R\otimes_KK'$ est réduit (ce qui aura lieu en
  particulier lorsque $K'$ est une extension séparable de $K$, $R$ étant
  composé direct de corps, extensions de $K$ (Bourbaki, \emph{Alg.},
  chap.~VIII, §~7, n\textsuperscript{o}~3, cor.~1 du th.~1)), alors $A'$ est
  un $A$-module de type fini.
  \item[(ii)] Si en particulier $R$ est une $K$-algèbre séparable, alors $A$
  est un anneau japonais.
\end{enumerate}
\end{proposition}

\begin{proof}
\emph{(i)} Posons pour simplifier $A_1=\widehat A$,
$A'_1=A'\otimes_AA_1$, $K'_1=K'\otimes_AA_1$. Comme $A_1$ est un
$A$-module plat (0\textsubscript{I}, 7.3.5), $A'_1$ s'identifie à un
sous-anneau de $K'_1$ et est évidemment entier sur $A_1$. D'autre part, si
l'on pose $K_1=K\otimes_AA_1$, on peut écrire
$K'_1=K'\otimes_KK_1$~; comme $A$ est intègre et que $A_1$ est un $A$-module
plat, tout élément $\neq0$ de $A$ est $A_1$-régulier
(0\textsubscript{I}, 6.3.4)~; comme $K=S^{-1}A$, où $S=A-\{0\}$, on a
$K_1=S^{-1}A_1$, et la remarque précédente prouve que $K_1$ s'identifie à un
sous-anneau de l'anneau total des fractions $R$ de $A_1$. Puisque tout
$K$-module est plat, $K'_1$ s'identifie à un sous-anneau de $K'\otimes_KR$,
qui par hypothèse est réduit et est une $R$-algèbre finie. Désignons par
$\mathfrak q_i$ ($1\leqslant i\leqslant n$) les idéaux premiers minimaux de
$A_1$, par $B_i$ l'anneau intègre $A_1/\mathfrak q_i$, par $L_i$ le corps des
fractions de $B_i$, de sorte que $R$ est composé direct des $L_i$, et
$K'\otimes_KR$ composé direct des $K'\otimes_KL_i$~; ces dernières
$K$-algèbres, étant réduites, sont des composées directes de corps,
extensions finies des $L_i$. Comme $B_i$ est un anneau local intègre
\emph{complet}, le théorème de Nagata (23.1.5) prouve que la fermeture
intégrale de $B_i$ dans $K'\otimes_KL_i$ est un $B_i$-module de type fini, et
\emph{a fortiori} un $A_1$-module de type fini~; comme tout élément de
$K'\otimes_KL_i$ qui est entier sur $A_1$ est aussi entier sur $B_i$ (étant
annulé par $\mathfrak q_i$), on en conclut finalement que la fermeture
intégrale de $A_1$ dans $K'\otimes_KR$ est un $A_1$-module de type fini. Mais
comme $A'_1$ est contenu dans cette fermeture intégrale et que $A_1$ est
noethérien, $A'_1$ est aussi un $A_1$-module de type fini. Enfin, comme $A_1$
est un $A$-module fidèlement plat (0\textsubscript{I}, 7.3.5), on en conclut
que $A'$ est un $A$-module de type fini (Bourbaki, \emph{Alg. comm.},
chap.~I, §~3, n\textsuperscript{o}~6, prop.~11).

\oldpage[0\textsubscript{IV-1}]{217}

\emph{(ii)} L'hypothèse entraîne que les $L_i$ sont des extensions séparables
de $K$, donc $K'\otimes_KL_i$ est réduit (Bourbaki, \emph{Alg.}, chap.~VIII,
§~7, n\textsuperscript{o}~3, th.~1) et l'on peut appliquer (i) à toute
extension finie $K'$ de $K$, ce qui démontre notre assertion.
\end{proof}

\subsection{Clôture intégrale d'un anneau local noethérien intègre}
\label{subsection:0.23.2-fr}

\begin{env}[23.2.1]
\label{0.23.2.1-fr}
Dans ce qui suit, pour tout anneau $A$, nous noterons $A_{\mathrm{red}}$ pour
simplifier, le quotient de $A$ par son nilradical (de sorte que si
$X=\operatorname{Spec}(A)$, on a
$X_{\mathrm{red}}=\operatorname{Spec}(A_{\mathrm{red}})$).

Nous dirons qu'un anneau \emph{local} $A$ est \emph{unibranche} si l'anneau
$A_{\mathrm{red}}$ est intègre et si la clôture intégrale de
$A_{\mathrm{red}}$ est un anneau \emph{local}, ce qui généralise la définition
donnée dans (III, 4.3.6). Nous dirons que $A$ est \emph{géométriquement
unibranche} s'il est unibranche et si le corps résiduel de l'anneau local,
clôture intégrale de $A_{\mathrm{red}}$, est une extension \emph{radicielle}
de celui de $A$. Il est clair qu'un anneau local intégralement clos est
géométriquement unibranche~; il résulte de (23.1.6) qu'un anneau local
noethérien intègre \emph{complet} est unibranche.
\end{env}

\begin{lemma}[23.2.2]
\label{0.23.2.2-fr}
\textnormal{(*)}\footnote{Dans la suite de ce numéro, nous utilisons des
notions et résultats exposés au chap.~IV, §§~2 et 5~; comme les résultats de
ce numéro ne sont pas utilisés avant le chap.~IV, §~6, cela n'entraîne pas de
cercle vicieux.} Soient $A$ un anneau intègre, $A'$ sa clôture intégrale~; pour
que le morphisme canonique
$f:\operatorname{Spec}(A')\to\operatorname{Spec}(A)$ soit radiciel, il faut et
il suffit que pour tout $\mathfrak p\in\operatorname{Spec}(A)$,
$A_{\mathfrak p}$ soit géométriquement unibranche~; lorsqu'il en est ainsi,
$f$ est un homéomorphisme.
\end{lemma}

\begin{proof}
En effet, pour tout $\mathfrak p\in\operatorname{Spec}(A)$, la clôture
intégrale de $A_{\mathfrak p}$ est $A'_{\mathfrak p}$ (Bourbaki,
\emph{Alg. comm.}, chap.~V, §~1, n\textsuperscript{o}~5, prop.~16), et tous
les idéaux premiers de $A'_{\mathfrak p}$ au-dessus de
$\mathfrak pA_{\mathfrak p}$ sont maximaux (\emph{loc. cit.}, §~2,
n\textsuperscript{o}~1, prop.~1). Dire que $f$ est injectif signifie donc que
pour tout $\mathfrak p\in\operatorname{Spec}(A)$, $A'_{\mathfrak p}$ est un
anneau \emph{local}, c'est-à-dire que les $A_{\mathfrak p}$ sont
unibranches. Dire que tout $A_{\mathfrak p}$ est géométriquement unibranche
signifie alors que $f$ est radiciel en vertu de (I, 3.5.8). Lorsqu'il en est
ainsi, $f$ est surjectif et fermé (II, 6.1.10), donc un homéomorphisme.
\end{proof}

\begin{lemma}[23.2.3]
\label{0.23.2.3-fr}
Soient $A$ un anneau intègre, $K$ son corps des fractions, $B$ un anneau
noethérien, $\varphi:A\to B$ un homomorphisme d'anneau faisant de $B$ un
$A$-module plat. Soit $K'$ une extension de $K$~; posons
$B_1=B_{\mathrm{red}}$, et soit $R$ l'anneau total des fractions de $B_1$.
Alors, pour toute sous-$A$-algèbre $A'$ de $K'$, l'homomorphisme canonique
\begin{equation}
\label{0.23.1.8.1-fr}
  (A'\otimes_AB)_{\mathrm{red}}
  \longrightarrow(K'\otimes_AR)_{\mathrm{red}}
  \tag{23.1.8.1}
\end{equation}
est injectif.
\end{lemma}

\begin{proof}
On peut en effet considérer l'homomorphisme canonique
$h:A'\otimes_AB\to K'\otimes_AR$ comme composé des homomorphismes suivants
\[
  A'\otimes_AB\xrightarrow{u}K'\otimes_AB
  \xrightarrow{v}K'\otimes_AB_1\xrightarrow{w}K'\otimes_AR.
\]
Comme $B$ est un $A$-module plat, $u$ est injectif~; de même, $K'$ étant un
$K$-module plat et $K$ un $A$-module plat, $K'$ est un $A$-module plat, et
$w$ est donc injectif puisque $B_1$ est réduit, donc l'homomorphisme
$B_1\to R$ injectif. Enfin, pour la même raison, si $\mathfrak N$ est

\oldpage[0\textsubscript{IV-1}]{218}

le nilradical de $B$, le noyau de $v$ est $K'\otimes_A\mathfrak N$, donc il est
contenu dans le nilradical de $K'\otimes_AB$~; on en conclut aussitôt que si
l'image par $h=w\circ v\circ u$ d'un élément $x\in A'\otimes_AB$ est
nilpotent, $x$ est nilpotent, ce qui prouve le lemme.
\end{proof}

\begin{proposition}[23.2.4]
\label{0.23.2.4-fr}
Soient $A$ un anneau intègre, $K$ son corps des fractions, $B$ un anneau
noethérien, $\mathfrak q_i$ ($1\leqslant i\leqslant n$) ses idéaux premiers
minimaux, $\varphi:A\to B$ un homomorphisme d'anneaux. On suppose que les
$B/\mathfrak q_i$ sont des anneaux japonais, et que $B$ soit un $A$-module
plat. Soit $K'$ une extension finie de $K$, et soit $(C_\lambda)$ la famille
filtrante croissante des sous-anneaux de $K'$ qui sont des $A$-algèbres
finies et admettent $K'$ pour corps des fractions~; la réunion $A'$ des
$C_\lambda$ est donc la fermeture intégrale de $A$ dans $K'$. Alors~:
\begin{enumerate}
  \item[(i)] Il existe un indice $\alpha$ tel que pour $\lambda\geqslant
  \alpha$, l'homomorphisme canonique
  \begin{equation}
  \label{0.23.2.4.1-fr}
    (C_\alpha\otimes_AB)_{\mathrm{red}}
    \longrightarrow(C_\lambda\otimes_AB)_{\mathrm{red}}
    \tag{23.2.4.1}
  \end{equation}
  soit bijectif.
  \item[(ii)] Si en outre $B$ est un $A$-module fidèlement plat, le morphisme
  $\operatorname{Spec}(A')\to\operatorname{Spec}(C_\alpha)$ est radiciel.
\end{enumerate}
\end{proposition}

\begin{proof}
\emph{(i)} Soit $B_1=B_{\mathrm{red}}$, et soit $R$ l'anneau total des
fractions de $B_1$~; comme $A$ est intègre et $B$ un $A$-module plat, tout
élément $\neq0$ de $A$ est $B$-régulier (0\textsubscript{I}, 6.3.4), donc
l'homomorphisme composé $A\xrightarrow{\varphi}B\to B_1$ se prolonge en un
homomorphisme $K\to R$~; on peut alors écrire
$K'\otimes_AR=(K'\otimes_AK)\otimes_KR$, et comme $K'\otimes_AK$
s'identifie à $K'$, on a $K'\otimes_AR=K'\otimes_KR$. Comme $R$ est composé
direct des corps des fractions $L_i$ des $B/\mathfrak q_i$,
$K'\otimes_AR$ est composé direct des $K'\otimes_KL_i$, et par suite
$(K'\otimes_AR)_{\mathrm{red}}$ est composé direct d'un nombre fini
d'extensions finies des $L_i$. En vertu de l'hypothèse, la fermeture
intégrale de $B/\mathfrak q_i$ dans une extension finie de $L_i$ est un
$(B/\mathfrak q_i)$-module de type fini, donc un $B$-module de type fini~;
on en conclut que la fermeture intégrale $B'$ de $B$ dans
$(K'\otimes_AR)_{\mathrm{red}}$ est un $B$-module de type fini. En vertu de
(23.2.3), les $(C_\lambda\otimes_AB)_{\mathrm{red}}$ s'identifient
canoniquement à des sous-anneaux de $(K'\otimes_AR)_{\mathrm{red}}$ qui sont
des $B$-algèbres finies, donc contenues dans $B'$. Comme $B$ est noethérien
et $B'$ un $B$-module de type fini, la famille filtrante des
$(C_\lambda\otimes_AB)_{\mathrm{red}}$ admet un \emph{plus grand élément}
$(C_\alpha\otimes_AB)_{\mathrm{red}}$, ce qui prouve (i).

\emph{(ii)} Comme $B$ est un $A$-module fidèlement plat, il suffit
(IV, 2.6.1, (v)) de montrer que le morphisme
$\operatorname{Spec}(A'\otimes_AB)\to
\operatorname{Spec}(C_\alpha\otimes_AB)$ est radiciel, ou, ce qui revient au
même (I, 5.1.6), que le morphisme
$\operatorname{Spec}(A'\otimes_AB)_{\mathrm{red}}\to
\operatorname{Spec}(C_\alpha\otimes_AB)_{\mathrm{red}}$ est radiciel~; mais
on a $A'\otimes_AB=\varinjlim_\lambda(C_\lambda\otimes_AB)$, donc
(IV, 5.13.2)
$(A'\otimes_AB)_{\mathrm{red}}=
\varinjlim_\lambda(C_\lambda\otimes_AB)_{\mathrm{red}}$, et la conclusion
résulte de (i).
\end{proof}

\begin{corollary}[23.2.5]
\label{0.23.2.5-fr}
Soient $A$ un anneau local noethérien intègre, $K$ son corps des fractions,
$K'$ une extension finie de $K$. Soit $(C_\lambda)$ la famille filtrante
croissante des sous-anneaux de $K'$ qui sont des $A$-algèbres finies (donc
des anneaux noethériens semi-locaux (Bourbaki, \emph{Alg. comm.}, chap.~IV,
§~2, n\textsuperscript{o}~5, cor.~3 de la prop.~9)) et admettent $K'$ pour
corps des fractions. Alors il existe $\alpha$ tel que l'homomorphisme
$(\widehat C_\alpha)_{\mathrm{red}}\to
(\widehat C_\lambda)_{\mathrm{red}}$ soit un isomorphisme pour
$\lambda\geqslant\alpha$, et si $A'$ est la fermeture intégrale de $A$ dans
$K'$, le morphisme
$\operatorname{Spec}(A')\to\operatorname{Spec}(C_\alpha)$ est radiciel
(cf. (23.2.2)).
\end{corollary}

\oldpage[0\textsubscript{IV-1}]{219}

\begin{proof}
On applique (23.2.4) en prenant $B=\widehat A$~; comme les
$B/\mathfrak q_i$ sont des anneaux locaux noethériens complets, ce sont des
anneaux japonais (23.1.5) et $B$ est un $A$-module fidèlement plat
(0\textsubscript{I}, 7.3.5)~; en outre on a
$C_\lambda\otimes_AB=\widehat C_\lambda$ (Bourbaki, \emph{Alg. comm.},
chap.~III, §~3, n\textsuperscript{o}~4, th.~3 et chap.~IV, §~2,
n\textsuperscript{o}~5, cor.~3 de la prop.~9).
\end{proof}

\begin{corollary}[23.2.6]
\label{0.23.2.6-fr}
Sous les hypothèses de (23.2.5), la fermeture intégrale $A'$ de $A$ dans
$K'$ est un anneau semi-local~; si $\mathfrak m$ est l'idéal maximal de $A$,
alors, pour tout idéal maximal $\mathfrak m'$ de $A'$,
$A'/\mathfrak m'$ est une extension finie de $A/\mathfrak m$.
\end{corollary}

\begin{proof}
Le fait que l'homomorphisme
$(\widehat C_\alpha)_{\mathrm{red}}\to
(\widehat C_\lambda)_{\mathrm{red}}$ soit bijectif pour
$\lambda\geqslant\alpha$ entraîne que le nombre des idéaux maximaux de
$C_\lambda$ est constant pour $\lambda\geqslant\alpha$ et que si
$\mathfrak m_\alpha$ est un idéal maximal de $C_\alpha$ et
$\mathfrak m_\lambda$ l'unique idéal maximal de $C_\lambda$ au-dessus de
$\mathfrak m_\alpha$, les corps $C_\alpha/\mathfrak m_\alpha$ et
$C_\lambda/\mathfrak m_\lambda$ sont canoniquement isomorphes. La conclusion
résulte de ce que $\mathfrak m'=\varinjlim\mathfrak m_\lambda$ et
$A'/\mathfrak m'=\varinjlim(C_\lambda/\mathfrak m_\lambda)$
(0\textsubscript{III}, 10.3.1.1.3).
\end{proof}

\begin{proposition}[23.2.7]
\label{0.23.2.7-fr}
\textnormal{(Y. Mori).} Soient $A$ un anneau local noethérien intègre, $K$ son
corps des fractions, $A'$ la clôture intégrale de $A$. Alors $A'$ est un
anneau de Krull (Bourbaki, \emph{Alg. comm.}, chap.~VII, §~1) semi-local~;
autrement dit, il existe une famille $(V_\lambda)$ d'anneaux de valuation
discrète ayant $K$ pour corps des fractions, tels que
$A'=\bigcap_\lambda V_\lambda$ et que tout $x\in K$ appartienne à tous les
$V_\lambda$ sauf pour un nombre fini d'entre eux. En outre, il existe un
sous-anneau $C$ de $K$ qui est une $A$-algèbre finie, et telle que les
$V_\lambda$ soient les clôtures intégrales des anneaux
$C_{\mathfrak p_\lambda}$, où $(\mathfrak p_\lambda)$ est la famille des
idéaux premiers de hauteur $1$ de l'anneau local $C$.
\end{proposition}

Prouvons d'abord le lemme suivant~:

\begin{lemma}[23.2.7.1]
\label{0.23.2.7.1-fr}
Soient $A$, $B$ deux anneaux, $\varphi:A\to B$ un homomorphisme faisant de $B$
un $A$-module fidèlement plat. On suppose que $A$ est intègre~; alors, si
$C=B_{\mathrm{red}}$, l'homomorphisme composé
$\psi:A\to B\to B_{\mathrm{red}}=C$ est injectif et se prolonge en un
homomorphisme injectif du corps des fractions $K$ de $A$ dans l'anneau total
des fractions $R$ de $C$~; en outre, si $C'$ est la fermeture intégrale de
$C$ dans $R$, $C'\cap K$ est la clôture intégrale de $A$.
\end{lemma}

\begin{proof}
Comme $A$ est intègre et $\varphi$ injectif, l'intersection de $\varphi(A)$
et du nilradical $\mathfrak N$ de $B$ est réduite à $0$, et comme
$C=B/\mathfrak N$, $\psi$ est injectif. Tout $a\neq0$ dans $A$ étant non
diviseur de zéro dans $A$ ne l'est pas non plus dans $B$ par platitude
(0\textsubscript{I}, 6.3.4)~; on en déduit que $a$ n'est pas non plus diviseur
de $0$ dans $C$, car si l'on avait $ax\in\mathfrak N$ pour un
$x\notin\mathfrak N$ dans $B$, on en tirerait $a^nx^n=0$ pour un entier $n$,
ce qui contredit ce qui précède puisque $x^n\neq0$. On peut donc prolonger
$\psi$ en un homomorphisme injectif de $K$ dans $R$. Pour prouver la dernière
assertion, notons qu'il est clair que la clôture intégrale $A'$ de $A$ est
contenue dans $C'\cap K$. Inversement, soit $x\in C'\cap K$~; $x$ est donc
entier sur $C$, et \emph{a fortiori} sur $B$, autrement dit $B[x]$ est une
$B$-algèbre finie~; d'ailleurs, $K$ s'identifie à un sous-anneau de
$B\otimes_AK$, et le sous-anneau $B[x]$ de $B\otimes_AK$ s'identifie à
$B\otimes_AA[x]$ par platitude~; on en conclut que $A[x]$ est un $A$-module
de type fini (Bourbaki, \emph{Alg. comm.}, chap.~I\textsuperscript{er}, §~3,
n\textsuperscript{o}~6, prop.~11), donc que $x\in A'$.
\end{proof}

Ce lemme étant établi, nous allons l'appliquer à l'injection canonique de $A$
dans son complété $\widehat A$, qui est un $A$-module fidèlement plat~;
$B=(\widehat A)_{\mathrm{red}}=\widehat A/\mathfrak N$ (où $\mathfrak N$ est
le nilradical de $\widehat A$) est un anneau local noethérien complet réduit,
dont l'anneau total des

\oldpage[0\textsubscript{IV-1}]{220}

fractions $R$ est donc composé direct d'un nombre fini de corps $L_i$~;
l'image canonique $B_i$ de $B$ dans $L_i$ est un anneau local noethérien
\emph{intègre et complet}, dont $L_i$ est le corps de fractions, et la
fermeture intégrale $B'$ de $B$ dans $R$ est \emph{composée directe} des
$B'_i$, où $B'_i$ est la clôture intégrale de $B_i$. Mais en vertu de
(23.1.5), $B'_i$ est une $B_i$-algèbre \emph{finie}, donc un anneau local
\emph{noethérien} intégralement clos. On sait (Bourbaki,
\emph{Alg. comm.}, chap.~VII, §~1, n\textsuperscript{o}~3, cor. du th.~2) que
$B'_i$ est un anneau de Krull. Or, pour tout $i$, on a un homomorphisme
$\varphi_i:K\to L_i$, qui est injectif, et par suite
$\varphi_i^{-1}(B'_i)$ est un anneau de Krull dans $K$. Mais comme
$A'=B'\cap K$ en vertu du lemme et que
$B'=\bigcap_i\varphi_i^{-1}(B'_i)$, $A'$ est intersection d'un nombre fini
d'anneaux de Krull et est par suite un anneau de Krull.

Pour prouver la dernière assertion de la proposition, notons qu'il existe une
$A$-algèbre finie $C\subset K$ telle que le morphisme
$\operatorname{Spec}(A')\to\operatorname{Spec}(C)$ soit radiciel et un
homéomorphisme (23.2.4)~; pour tout idéal premier
$\mathfrak p'\in\operatorname{Spec}(A')$, $\mathfrak p'$ est donc le seul
idéal premier de $A'$ au-dessus de $\mathfrak p=\mathfrak p'\cap C$ et
$A'_{\mathfrak p'}$ la clôture intégrale de $C_{\mathfrak p}$~; d'autre part
le fait que $\operatorname{Spec}(A')\to\operatorname{Spec}(C)$ est un
homéomorphisme entraîne que l'application
$\mathfrak p'\mapsto\mathfrak p'\cap C$ est une bijection de l'ensemble des
idéaux premiers de hauteur $1$ de $A'$ sur l'ensemble des idéaux premiers de
hauteur $1$ de $C$~; d'où la conclusion, compte tenu de Bourbaki,
\emph{Alg. comm.}, chap.~VII, §~1, n\textsuperscript{o}~6, th.~4.

\begin{remarks}[23.2.8]
\label{0.23.2.8-fr}
\begin{enumerate}
  \item[(i)] On aura soin de noter, dans l'application de (23.2.6) et
  (23.2.7), que l'anneau $A'$ \emph{n'est pas nécessairement un anneau
  noethérien}, comme le montre un exemple de Nagata avec $\dim(A)=3$ [30].

  \item[(ii)] Les conclusions de (23.2.7) sont encore valables si $A'$ est
  la fermeture intégrale de $A$ dans une extension finie $K'$ de $K$~; il
  suffit en effet de considérer une $A$-algèbre finie $B$ dont $K'$ est le
  corps des fractions~; $B$ est un anneau semi-local noethérien, donc
  intersection d'un nombre fini d'anneaux locaux noethériens
  $B_{\mathfrak m_i}$ ($\mathfrak m_i$ idéaux maximaux de $B$), et sa clôture
  intégrale (qui est égale à $A'$) est intersection des
  $A'_{\mathfrak m_i}$ (Bourbaki, \emph{Alg. comm.}, chap.~II, §~3,
  n\textsuperscript{o}~3, cor.~4 du th.~1) qui sont les clôtures intégrales
  des $B_{\mathfrak m_i}$, donc des anneaux de Krull par (23.1.13)~; par
  suite $A'$ est un anneau de Krull.

  \item[(iii)] On peut montrer ([30], 3.3.10) que pour tout anneau noethérien
  intègre $A$ (non nécessairement local), la clôture intégrale de $A$ est un
  anneau de Krull.
\end{enumerate}
\end{remarks}

\begin{corollary}[23.2.9]
\label{0.23.2.9-fr}
Soient $A$ un anneau noethérien intègre, $A'$ sa clôture intégrale. Supposons
que pour tout anneau $C$ tel que $A\subset C\subset A'$ et tel que $C$ soit
une $A$-algèbre finie, et pour tout idéal premier $\mathfrak q$ de hauteur $1$
dans $C$, $\mathfrak q\cap A$ soit un idéal premier de hauteur $1$ dans $A$.
Soit $P$ l'ensemble des idéaux premiers de hauteur $1$ dans $A$~; on a alors
\[
  A'=\bigcap_{\mathfrak p\in P}A'_{\mathfrak p}.
\]
\end{corollary}

\begin{proof}
Comme $A'$ est un $A$-module sans torsion, on a
$A'=\bigcap A'_{\mathfrak m}$, où $\mathfrak m$ parcourt l'ensemble des idéaux
maximaux de $A$ (Bourbaki, \emph{Alg. comm.}, chap.~II, §~3,
n\textsuperscript{o}~3, cor.~4 du th.~1)~;

\oldpage[0\textsubscript{IV-1}]{221}

il suffira donc de prouver que $A'_{\mathfrak m}$ est l'intersection des
$A'_{\mathfrak p}$ pour les idéaux premiers $\mathfrak p\subset\mathfrak m$
de hauteur $1$. Comme $A'_{\mathfrak m}$ est la clôture intégrale de
$A_{\mathfrak m}$, on voit qu'on peut se borner à démontrer le corollaire pour
$A_{\mathfrak m}$. Or, il y a alors, en vertu de (23.2.7), une
$A_{\mathfrak m}$-algèbre finie $B$ contenue dans $A'_{\mathfrak m}$, telle
que $A'_{\mathfrak m}$ soit l'intersection des clôtures intégrales des anneaux
$B_{\mathfrak q}$ où $\mathfrak q$ parcourt l'ensemble des idéaux premiers de
hauteur $1$ de $B$. Mais $B$ est de la forme $C_{\mathfrak m}$, où $C$ est
une $A$-algèbre finie~; donc, pour tout idéal premier $\mathfrak q$ de hauteur
$1$ dans $B$, $\mathfrak q\cap A$ est par hypothèse un idéal premier de
hauteur $1$ dans $A$~; comme la clôture intégrale de $B_{\mathfrak q}$
contient celle de $A_{\mathfrak q\cap A}$ et que cette dernière contient
$A'_{\mathfrak m}$, il en résulte bien que $A'_{\mathfrak m}$ est
l'intersection des $A'_{\mathfrak p}$ pour $\mathfrak p\in P$ et
$\mathfrak p\subset\mathfrak m$, ce qui achève la démonstration.
\end{proof}

\begin{remarks}[23.2.10]
\label{0.23.2.10-fr}
\begin{enumerate}
  \item[(i)] Nous verrons (IV, 5.6.3 et 5.6.10) que l'hypothèse de (23.2.9)
  est vérifiée lorsque $A$ est \emph{universellement caténaire}, en particulier
  (IV, 5.10.17 et 5.11.2) lorsque $A$ est \emph{quotient d'un anneau
  régulier}~; enfin, elle est évidemment vérifiée si
  $\operatorname{Spec}(A')\to\operatorname{Spec}(A)$ est un homéomorphisme,
  en particulier si ce morphisme est radiciel (23.2.2).

  \item[(ii)] La conclusion de (23.2.9) n'est plus nécessairement exacte
  lorsqu'on omet l'hypothèse sur les sous-$A$-algèbres finies de $A'$. En
  effet, on peut définir un anneau local noethérien intègre $A$, de dimension
  $2$, dont la clôture intégrale $A'$ est une $A$-algèbre finie, mais tel qu'il
  y ait un idéal premier $\mathfrak p'$ de $A'$ de hauteur $1$ au-dessus de
  l'idéal maximal $\mathfrak m$ (de hauteur $2$) de $A$, et tel en outre que
  pour tout idéal premier $\mathfrak p$ de hauteur $1$ dans $A$,
  $A_{\mathfrak p}$ soit intégralement clos (IV, 5.6.11)~; l'intersection de
  ces anneaux $A_{\mathfrak p}$ ne peut alors être égale à $A'$ puisque $A'$
  est un anneau de Krull (Bourbaki, \emph{Alg. comm.}, chap.~VII, §~1,
  n\textsuperscript{o}~5, prop.~9).
\end{enumerate}
\end{remarks}

\begin{flushright}
\emph{(A suivre.)}
\end{flushright}
